\documentclass[11pt]{article}
\usepackage[margin=1in]{geometry}
\usepackage{amsmath,amssymb,amsthm}
\usepackage{booktabs}
\usepackage{enumitem}
\usepackage[hidelinks]{hyperref}

\newtheorem{definition}{Definition}[section]
\newtheorem{proposition}[definition]{Proposition}
\newtheorem{remark}[definition]{Remark}

\title{Spherepop Experiment Spec: Multi-Timescale Continuation}
\author{Flyxion}
\date{2026}

\begin{document}
\maketitle

\section{Purpose}

This experiment family operationalizes the bridge between \emph{The Geometry of Boredom}
and Spherepop by modeling boredom as a scope-selection problem over multiple unresolved
histories with heterogeneous temporal horizons.

\begin{remark}[Disambiguation from Appendix B Plan B]
This experiment suite is unrelated to the repository's pending Appendix B poset-integration
rewrite. It does not change Spherepop's four primitives or claim that full Appendix B
semantics is already implemented. It introduces a separate scheduling formalization for
multi-timescale continuation using existing vocabulary.
\end{remark}

\section{Core Formalization}

\begin{remark}[Grounding and status]
This document treats EMCG-based scope scheduling as a \emph{new proposed extension}, not as
an already-established theorem in the existing Spherepop specification. A related but distinct
saturation notion already exists in this corpus via the Asymptotic Saturation Theorem in
\texttt{textbook/The\_Ecology\_of\_Distinctions.tex} (\S ``Asymptotic Saturation Without
Exhaustion''), which concerns reintegration dynamics rather than task-portfolio scheduling.
The present experiment asks whether an analogous local-vs-field saturation distinction is
useful in practice at the workflow level.
\end{remark}

\begin{definition}[Scope portfolio]
A portfolio is
\[
\mathcal{T}=\{(S_i,\tau_i)\}_{i=1}^n,\qquad S_i=(O_i,\tau_i),\qquad
\tau_1\ll\tau_2\ll\cdots\ll\tau_n,
\]
where each unresolved scope $S_i$ has local option space $O_i$ and horizon $\tau_i$.
\end{definition}

\begin{definition}[Local and field saturation]
With operator $F_t=C(H_t)$,
\[
\text{LocalSat}(i,t):\ \max_{a\in O_i}\mathrm{EMCG}_i(a\mid F_t)\approx0,
\]
\[
\text{FieldSat}(t):\ \max_i\max_{a\in O_i}\mathrm{EMCG}_i(a\mid F_t)\approx0.
\]
\end{definition}

\begin{proposition}[Scope-switch recovery]
$\text{LocalSat}(i,t)$ does not imply $\text{FieldSat}(t)$. If another scope $j$ satisfies
$\max_{a\in O_j}\mathrm{EMCG}_j(a\mid F_t)>0$, switching scopes restores progress.
\end{proposition}

\begin{proposition}[Operator-drift recovery]
A scope can transition from
$\mathrm{EMCG}_i(a\mid F_t)\approx0$ to
$\mathrm{EMCG}_i(a\mid F_{t+\Delta t})>0$ without changing $O_i$, due to intervening work on
other scopes that modifies $F$.
\end{proposition}

\begin{proposition}[BIND-as-progress]
Long-horizon scopes can accumulate useful structure through BIND/REFUSE/COLLAPSE while
remaining unresolved; progress does not require immediate POP.
\end{proposition}

\section{Scheduling Policies Under Comparison}

\begin{enumerate}[leftmargin=*]
  \item \textbf{Novelty-only}: prefer longest-idle scope.
  \item \textbf{Shortest-task-first}: always choose minimal horizon.
  \item \textbf{Max-EMCG}: greedily maximize immediate expected marginal compression gain.
  \item \textbf{EMCG/cost + anti-starvation}: maximize gain-per-cost with forced periodic
  updates to slow scopes.
\end{enumerate}

\section{Event Semantics in This Experiment Family}

The run model interprets Spherepop operations as:
\[
\begin{aligned}
\mathrm{POP}&: \text{resolve ready local scope},\\
\mathrm{REFUSE}&: \text{remove continuation branch},\\
\mathrm{BIND}&: \text{tighten continuation constraints},\\
\mathrm{COLLAPSE}&: \text{merge distinctions no longer worth explicit tracking}.\\
\end{aligned}
\]

\section{Metrics}

\begin{itemize}[leftmargin=*]
  \item Time-to-first-local-saturation
  \item Time-in-field-saturation
  \item Mean recovery latency after scope switch
  \item Long-horizon starvation rate
  \item Compression-progress stability across horizons
\end{itemize}

\section{Deliverables and Outputs}

\begin{itemize}[leftmargin=*]
  \item Executable policy comparison script: \texttt{run.py}
  \item Results table template: \texttt{results-template.csv}
  \item Policy summaries including per-scope event counts and saturation statistics
\end{itemize}

\section{Application Mapping}

\begin{description}[leftmargin=2em,style=nextline]
  \item[Personal workflow] Maintain active scopes at minute/day/week/year horizons.
  \item[Education] Alternate short and long-horizon work instead of novelty churn.
  \item[Research/programming] Keep deep threads unresolved while adding regular BIND-like constraints.
  \item[AI agents] Use multi-clock schedulers with anti-starvation guarantees.
\end{description}

\end{document}
